\enteteExam{NFP 107 -- Systèmes de gestion de bases de données - Durée 2h30}%
{Première session  -- 3 février 2021}%
{Effectué à distance}%


On veut implanter une application de réservation de places en crèche. Voici le schéma proposé\,:

\begin{itemize}
  \item Crèche (\textbf{idCrèche}, nom, ville, capacité)
   \item Personne (\textbf{idPersonne}, prénom, nom, néEn, ville, idPère, idMère)
  \item Demande (\textbf{idDemande}, idEnfant, idCrèche, année)
  \item Inscription (\textbf{idEnfant, année}, idCrèche, idPuéricultrice)
   
\end{itemize}

\section{Modélisation (8 points)}


\begin{enumerate}
  \item Etudiez soigneusement le schéma et identifiez toutes les clés étrangères
      qui font référence à la table \texttt{Personne} (1 pt)
     
     
\begin{solution}
Il y en a beaucoup\,: idPère, idMère, idEnfant et idPuericultrice.
\end{solution}
 
  \item Donnez les commandes SQL de création des tables suivantes\,: 
  Crèche, Personne et Inscription. Note\,: on peut ne pas connaître
  le père ou la mère d'un enfant. (2 pts)
  
  
\begin{solution}
Crèche\: trivial.

\smallskip

\begin{small}
\begin{minted}{sql}
create table Personne 
   (idPersonne int not null,
    prénom varchar not null,
    nom varchar not null,
    néEn int not null,
    ville varchar not null,
    idPère int ,
    idMère int ,
    primary key (idPersonne),
    foreign key (idPère) reference Personne(idPersonne),
    foreign key (idMère) reference Personne(idPersonne)
   )

create table Inscription 
   (idEnfant int not null,
    année int not null,
    idCrèche int not null,
    idPuericultrice int not null,
    primary key (idEnfant, année),
    foreign key (idEnfant) reference Personne(idPersonne),
    foreign key (idCrèche) reference Crèche(idCrèche),
    foreign key (idPuericultrice) reference Personne(idPersonne)
   )
\end{minted}
\
\end{small}
\end{solution}
  
  \item Voyez-vous une table résultant d'une réification\,? Une table résultant d'une entité faible\,? 
     Lesquelles\,? (1 pt)
     
  \item Donnez le schéma Entité/association correspondant à ce schéma relationnel (2 pts)
  
  \item Un enfant peut-il être inscrit à plusieurs crèches la même année\,; un enfant peut-il demander
      à s'inscrire dans plusieurs crèches la même année\,? (1 pt)
      
\begin{solution}
    Non, pas d'inscription d'un enfant la même année car la clé est constituée de ces deux attributs.
\end{solution}


    \item Si on ajoutait la dépendance fonctionnelle \texttt{idPuéricultrice $\to$ idCrèche},
        comment faudrait-il modifier le schéma pour qu'il reste en 3FN\,? (1 pts)
      
  
\end{enumerate}

\section{Requêtes SQL (8 points)}

Les requêtes sont à exprimer en SQL. 

\begin{enumerate}

  
  \item Quels sont les enfants inscrits dans une crèche à Paris (donnez le nom de l'enfant
      et le nom de la crèche)\,?

\begin{solution}
\begin{small}
\begin{minted}{sql}
select p.nom as 'nomEnfant', c.nom as 'nomCrèche'
from Crèche as c, Inscription as i, Personne as p
where c.idCrèche = i.idCrèche
and i.idEnfant=p.idPersonne
and c.ville='Paris'
\end{minted}
\end{small}
\end{solution}
  \item Quels sont les enfants demandant une crèche dans la ville où ils habitent
  (donner le nom, prénom et la ville)\,?
  
\begin{solution}
\begin{small}
\begin{minted}{sql}
select p.prénom, p.nom, p.ville
from Crèche as c, Demande as d, Personne as p
where c.idCrèche = d.idCrèche
and d.idEnfant=p.idPersonne
and c.ville=p.ville
\end{minted}
\end{small}
\end{solution}

  \item Pour chaque enfant donnez le prénom de l'enfant, le nom de son père, de sa mère
  et de sa puéricultrice.

\begin{solution}
\begin{small}
\begin{minted}{sql}
select enfant.nom as 'nomEnfant', père.nom as 'nomPère',
       mère.nom as 'nomMère', puericultrice.nom as 'nomPuericultrice'
from Inscription as i, Personne as enfant, Personne as père, 
        Personne as mère, Personne as puericultrice
where i.idEnfant = enfant.idPersonne
and   enfant.idPère = père.idPersonne
and   enfant.idMère = mère.idPersonne
and i.idPuericultrice = puericultrice.idPersonne
\end{minted}
\end{small}
\end{solution}

 \item Je veux donner la liste de tous les enfants nés en 2019, avec le nom
     de leur puericultrice s'ils sont inscrits à une crèche, et \texttt{null}
     sinon. Donnez la requête SQL.
     
\begin{solution}
 \begin{small}
\begin{minted}{sql}
select e.nom as 'nomEnfant',  p.nom as 'nomPuericultrice'
from  Personne as e left join 
        (Inscription as i join Personne as p on i.idPuericultrice=p.idPersonne)
       on  i.idEnfant = e.idPersonne
where néEn=2019
\end{minted}
\end{small}
\end{solution}
 
  \item Indiquez les demandes en cours\,: celles pour lesquelles il existe une 
  demande de l'enfant à une crèche, mais pas d'inscription. Donnez simplement
  l'identifiant de l'enfant et l'identifiant de la crèche.
 
\begin{solution}
 \begin{small}
\begin{minted}{sql}
select d.idEnfant, c.idCrèche
from  Demande as d
where not exists (select * from Inscription as i
                  where d.idEnfant = i.idEnfant
                  and d.idCrèche = i.idCrèche)
\end{minted}
\end{small}
\end{solution}
 
  \item Donnez la liste des enfants née en 2018 et les crèches de leur ville
     qu'ils n'ont pas demandées (nom de l'enfant, nom de la crèche).

\begin{solution}
\begin{small}
\begin{minted}{sql}
select p.nom as 'nomEnfant', c.nom as 'nomCrèche'
from  Personne as p, Crèche as c
where not exists (select * from Demande as d
                  where d.idEnfant = i.idEnfant
                  and d.idCrèche = i.idCrèche)
\end{minted}
\end{small}
\end{solution}
 
  \item Quelles sont les crèches n'ayant pas rempli toutes leurs places
  (autrement dit  le nombre d'inscriptions est inférieur à la capacité)?

\begin{solution}
\begin{small}
\begin{minted}{sql}
select c.*
from  Crèche as c, Inscription as i
where c.idCrèche = i.idCrèche
group by c.idCrèche
having count(*) < c.capacité
\end{minted}
\end{small}
\end{solution}
   Aide\,: faites une première requête affichant le nombre d'inscriptions
      et  la capacité. Puis ajoutez la clause indiquant que le premier doit
      être inférieur au second.
\end{enumerate}
 

\section{Algèbre (2 points)}

Donnez l'expression algébrique pour la requête 5 de la section précédente.


\begin{solution}
$$\pi_{idEnfant,idCreche} (Demande) - \pi_{idEnfant,idCreche} (Cr\acute{e}che)$$
\end{solution}

\section{Transactions (4 points)}

On implante une procédure d'inscription qui effectue les étapes suivantes\,:

\begin{enumerate}
  \item On parcourt les demandes des enfants qui ne sont pas encore inscrits
  \item Pour chaque demande on vérifie qu'il reste de la place dans la crèche demandée
  \item Si oui, on crée le nuplet dans la table \texttt{Inscription}
\end{enumerate}

\begin{minted}{bash}
Procédure inscription()
begin
  demandes_en_cours = (résultat requête 5)
  for each demande in demandes_en_cours do
    places_restantes = (résultat requête 6)
    si (places_restantes > 0)
      Insérer inscription
    # Emplacement 1
  end for
  # Emplacement 2
end
\end{minted}
\paragraph{Questions}

\begin{enumerate}
  \item Où faut-il placer un \texttt{commit} dans la procédure\,? Emplacement 1 ou 2. Expliquez
votre choix.
  \item Est-ce que l'insertion d'une inscription change le contenu du curseur sur les demandes
      en cours\,?  Justifiez votre réponse.
\end{enumerate}
\end{document}

